% Data-plot skeleton (pgfplots). Copy next to plotfig.sty, rename, replace the data.
% Pick the example in assets/examples/plots/ closest to the chart you need and start from that instead when
% the chart type is not a plain line plot: grouped-bars, dumbbell, pareto, heatmap, stacked-100, donuts, radar,
% curves-bands (two panels with confidence bands); flows.tex is generated by scripts/flows_from_csv.py.
% Build:  python <skill>/scripts/build_figure.py <name>.tex
% Include: \includegraphics{<name>.pdf} with no width option; set width= and height= here instead.
%
% Revision contract: data, series names, axis ranges and printed values change only when the user says so; a
% requested change to one element (a legend, a reference line, a label) leaves every other line as it is.
% Keep the `paper` style, the cycle list and the fonts; a figure that needs another look is a new template.
\documentclass[10pt,tikz,border=2pt]{standalone}
\usepackage{plotfig}
\begin{document}
\begin{tikzpicture}
\begin{axis}[paper, width=6cm, height=4cm, y label top,
  xlabel={training examples}, ylabel={accuracy (\%)},
  xmin=0, xmax=5, ymin=20, ymax=90, ytick={20,40,60,80}]  % direct labels hang outside the axis (clip=false); the page grows to fit them
% log axis: `xmode=log, log ticks x` (xmode itself must be written here, not inside a style); ticks then read
% 10, 100, 1000 in the text font. Zero cannot be shown on a log axis.
% before choosing the chart at all, read references/principles.md (which chart answers which question, and why)
\addplot[ref, forget plot] coordinates {(0,86) (5,86)} node[dl, text=black!55] {human};
\addplot coordinates {(0,42) (1,55) (2,66) (3,74) (4,79) (5,82)} node[dl, text=kept] {Ours};
\addplot coordinates {(0,38) (1,47) (2,56) (3,63) (4,68) (5,71)} node[dl, text=addctext] {System A};
% colours and solid marks come from the cycle list "paper": blue, orange, teal, violet, gold, grey. When the
% figure is "ours against the field", add cycle list name=paper accent to the axis (blue, then three greys).
% Orange text at this size is addctext, not addc.
% when two lines end close together: node[dl, yshift=3pt] and node[dl, yshift=-3pt]
% one annotation of the finding the text discusses:
% \node[note, anchor=north] at (axis cs:3,35) {A saturates\\ after 3k};  \draw[notearrow] (axis cs:3,36) -- (axis cs:3,60);
\end{axis}
\end{tikzpicture}
\end{document}
